\chapter{Arithmetic degrees and maximal slopes}

\begin{definition}[Hermitian bundles and normalized degree]
\label{def:hermitian-degree}
\source{bost-specht-manuscript:page-0002,bost-specht-manuscript:page-0003}\notready
A Hermitian bundle $\overline A$ is a finite projective $\Ocal$-module $A$
with a positive-definite Hermitian metric on each complex archimedean fiber,
invariant under conjugation, and a positive-definite real metric on each real
fiber. For a rank-one bundle $\overline L$ and $0\ne\ell\in L$, set
\[
 \degn(\overline L)=\frac{1}{D}\left(
 \log|L/\Ocal\ell|-\sum_v w_v\log\|\ell\|_v\right).
\]
The product formula makes this independent of $\ell$. For positive rank,
$\degn(\overline A)=\degn(\det\overline A)$ and
$\slope(\overline A)=\degn(\overline A)/\rk A$; set $\degn(0)=0$.
A bundle is semistable if every nonzero submodule has slope at most its own
slope, and strictly stable if every nonzero proper saturated submodule has
strictly smaller slope. It is normalized if its slope is zero.
\end{definition}

\begin{lemma}[Determinant identities and contraction]
\label{lem:degree-calculus}\uses{def:hermitian-degree}
\source{bost-specht-manuscript:page-0003}\notready
Degrees are additive in an exact sequence with induced and quotient
metrics. For positive-rank bundles $A,B$ and a line bundle $L$,
\[
 \slope(A\otimes B)=\slope(A)+\slope(B),\qquad
 \slope(A^\vee)=-\slope(A),\qquad
 \slope(L\otimes A)=\degn(L)+\slope(A).
\]
If $f:A\to B$ is an injective map of equal-rank bundles with
$\|f\|_v\leq1$ at every archimedean place, then
$\degn(A)\leq\degn(B)$. Replacing a submodule by its saturation cannot
decrease degree. Multiplying all norms by $e^t$ subtracts $t$ from every
positive-rank submodule's slope.
\end{lemma}
\begin{proof}
For an exact sequence, the determinant isomorphism
$\det A\otimes\det C\simeq\det B$ is an isometry: at each place choose
an orthonormal basis of the subspace and lift an orthonormal basis of the
quotient to its orthogonal complement. Local bases at finite primes show
that the corresponding determinant ideals multiply. Substituting these
two facts in the definition gives additivity. In orthonormal bases the
determinant identities
$\det(A\otimes B)\simeq(\det A)^{\otimes\rk B}\otimes
(\det B)^{\otimes\rk A}$ and $\det(A^\vee)\simeq(\det A)^\vee$
are isometries. Their local determinant ideals obey the same identities,
which proves all three slope formulas.

For an equal-rank injection, identify the generic fibers by $f$. The
finite-index inclusion of determinant modules contributes
$D^{-1}\log|B/f(A)|\geq0$ to $\degn(B)-\degn(A)$. The metric
contribution is
$-D^{-1}\sum_v w_v\log\|\det f\|_v\geq0$, since each singular value
of $f$ is at most one. A submodule and its saturation have identical
generic fibers and metrics, so this also proves the saturation claim.
Finally, a rank-$b$ determinant norm is multiplied by $e^{bt}$ under
uniform scaling, so its normalized degree decreases by $bt$.
\end{proof}

\begin{lemma}[Restriction of scalars and covolume]
\label{lem:restriction-covolume}\uses{def:hermitian-degree}
\source{bost-specht-manuscript:page-0011}\notready
Give $A\otimes_{\mathbb Z}\mathbb R$ the Euclidean norm
$\|x\|^2=\sum_v w_v\|x_v\|_v^2$. For $a=\rk_\Ocal A$,
\[
 \covol_\mathbb Z(A)=\Delta_K^{a/2}\exp(-D\degn(\overline A)).
\]
\end{lemma}
\begin{proof}
For $A=\Ocal^a$ with the standard metrics, the Gram determinant of an
integral basis of $\Ocal$ for this real trace metric is $\Delta_K$.
Indeed, the matrix of embeddings gives the trace discriminant; replacing
a conjugate pair by real and imaginary parts and weighting both by $2$
preserves its absolute squared determinant. Orthogonal direct sum gives
the standard covolume $\Delta_K^{a/2}$. If the fiber metrics have Gram
matrices $G_v$, covolume is multiplied by
$\prod_v(\det G_v)^{w_v/2}$, while the definition of degree multiplies
$\exp(-D\degn A)$ by precisely that quantity. This proves the formula
for a free module with arbitrary metrics. For general $A$, choose a
generic basis lying in $A$ and its finite-index free span $A_0$.
Then $\covol(A_0)=[A:A_0]\covol(A)$ and
$D\degn(A)=D\degn(A_0)+\log[A:A_0]$; cancellation proves the claim.
\end{proof}

\begin{lemma}[Finiteness of submodules of bounded degree below]
\label{lem:degree-finiteness}\uses{lem:restriction-covolume}
\source{bost-specht-manuscript:page-0012}\notready
For a fixed Hermitian bundle $A$, a rank $b$ and a real $c$, only finitely
many saturated rank-$b$ submodules $B\subseteq A$ satisfy $\degn(B)\geq c$.
\end{lemma}
\begin{proof}
Saturation over $\Ocal$ implies saturation over $\mathbb Z$: if a
nonzero integer multiple of $x\in A$ lies in $B$, torsion-freeness of
$A/B$ gives $x\in B$. A saturated rank-$Db$ sublattice of the underlying
$\mathbb Z$-lattice is determined by its primitive exterior vector in
$\bigwedge^{Db}_\mathbb Z A$, up to sign. The norm of this vector is its
covolume. By the preceding formula that norm is at most
$\Delta_K^{b/2}e^{-Dc}$. A bounded ball meets a discrete lattice in a
finite set, and hence only finitely many such primitive exterior vectors
occur. Their rational spans, and therefore their saturated sublattices,
determine the original submodules.
\end{proof}

\begin{lemma}[Finite-dimensional convex interval slopes]
\label{lem:subspace-convex-slopes}
\source{chen-jeannin-harder-narasimhan-games-v1}\notready
Let $V$ be a finite-dimensional vector space and let $\delta$ be a
real-valued function on its lattice of subspaces. For $X\subsetneq Y$ put
\[
 r(X,Y)=\dim Y-\dim X,\qquad d(X,Y)=\delta(Y)-\delta(X),\qquad
 \mu(X,Y)=\frac{d(X,Y)}{r(X,Y)}.
\]
Both $r$ and $d$ are additive along strictly increasing chains, and $r$
is strictly positive. For $X\subsetneq Y\subsetneq Z$, the three slopes
$\mu(X,Y),\mu(X,Z),\mu(Y,Z)$ are either all equal or strictly ordered
in the displayed order or its reverse. If
\[
 \delta(X)+\delta(Y)\leq\delta(X\cap Y)+\delta(X+Y),
\]
then for $X\not\subseteq Y$,
$\mu(X\cap Y,X)\leq\mu(Y,X+Y)$.
Every strictly increasing or strictly decreasing chain has at most
$\dim V$ strict steps.
\end{lemma}
\begin{proof}
Strict subspace inclusions strictly increase dimension. Differences of
dimensions and of $\delta$ telescope on a chain, giving additivity and
positivity. Thus $\mu(X,Z)$ is the weighted average of $\mu(X,Y)$ and
$\mu(Y,Z)$ with positive weights $r(X,Y),r(Y,Z)$; the three asserted
possibilities follow by subtracting either endpoint from this average.
The dimension identity
$\dim X+\dim Y=\dim(X\cap Y)+\dim(X+Y)$ makes the two denominators
in the convexity assertion equal and positive. Supermodularity of
$\delta$ gives the numerator inequality, so division proves it.
Finally, dimensions are integers between zero and $\dim V$; they
increase, respectively decrease, by at least one at every strict step.
\end{proof}

\begin{remark}
The slope-like part is the positive-rank specialization of the additive
rank-and-degree criterion in \cite[Section 4.3]{chen-jeannin}. Strict
dimension growth excludes its zero-rank case. Convexity additionally
uses supermodularity of the subspace degree; no modularity of that
degree is asserted.
\end{remark}

\begin{lemma}[The Hermitian subspace degree is supermodular]
\label{lem:hermitian-convex-slopes}
\uses{lem:degree-calculus,lem:subspace-convex-slopes}\notready
For a Hermitian bundle $A$ put $\delta(X)=\degn(A\cap X)$ on the
$K$-subspaces of $A_K$, using the induced metrics. Then $\delta$ is
supermodular. If $X\subsetneq Y$, the interval slope in the preceding
lemma equals the slope of $(A\cap Y)/(A\cap X)$ with its quotient
metric. In particular $\mu(0,A_K)=\slope(A)$, and the inequalities
$\mu(0,X)\leq\mu(0,A_K)$ for $X\ne0$ are exactly semistability.
\end{lemma}
\begin{proof}
Write $U=A\cap X$, $V=A\cap Y$, $I=U\cap V$ and $T=U+V$.
The map $U/I\to T/V$ is contracting: distance to $V$ is no greater
than distance to $I$. By degree additivity and contraction,
$\degn(U)+\degn(V)\leq\degn(I)+\degn(T)$. Saturating $T$ gives
$A\cap(X+Y)$ and can only increase degree. Also
$I=A\cap(X\cap Y)$, proving supermodularity. For nested subspaces,
degree additivity identifies the quotient's degree with
$\delta(Y)-\delta(X)$ and its rank with $\dim Y-\dim X$.
Every saturated submodule is $A\cap X$ for its generic fiber $X$,
and saturation increases degree without changing rank; hence checking
these subspaces is equivalent to checking all submodules.
\end{proof}

\begin{lemma}[Harder--Narasimhan filtration]
\label{lem:hn-filtration}
\uses{def:hermitian-degree,lem:degree-calculus,lem:degree-finiteness,lem:hermitian-convex-slopes}
\source{bost-specht-manuscript:page-0003}\notready
Every nonzero Hermitian bundle $A$ has a unique saturated filtration with
semistable quotient bundles and strictly decreasing slopes. The largest
slope of a nonzero submodule is attained, equals the first quotient's
slope, and is denoted by $\maxslope(A)$. The first filtration term is
the unique submodule of largest rank among those of slope $\maxslope(A)$.
Every isometric automorphism preserves the filtration.
\end{lemma}
\begin{proof}
For each possible positive rank, degree finiteness above any one attained
degree shows that the supremum of degrees is a maximum. There are finitely
many ranks, so maximal slope is attained. Choose a saturated submodule
$M$ of maximal slope $\mu$, with greatest rank among these.
For saturated submodules $U,V$, let $I=U\cap V$ and let $T$ be the
saturation of $U+V$. Supermodularity of the Hermitian subspace degree gives
\[
 \degn(U)+\degn(V)\leq\degn(I)+\degn(T).
\]
The right-hand side is at most
$\mu(\rk I+\rk T)=\mu(\rk U+\rk V)$. If $U,V$ have slope $\mu$,
equality follows throughout, including when $I=0$ using $\degn(0)=0$.
Thus $T$ also has slope $\mu$. Maximality of the rank
of $M$ forces every maximal-slope submodule to lie in $M$, proving its
uniqueness. It is semistable by the definition of $\mu$.

For a positive-rank submodule $N\subseteq A/M$, let $P\subseteq A$ be
its inverse image. The quotient metric on $P/M$ is the induced metric
on $N$: orthogonal projection of a vector in $P$ off $M$ stays in $P$.
Thus $\degn(P)=\degn(M)+\degn(N)$. If $\slope(N)>\mu$, the slope
of $P$ exceeds $\mu$; if $\slope(N)=\mu$, then $P$ has slope $\mu$
and greater rank. Both are impossible. Repeating on $A/M$ constructs
the filtration, with strictly decreasing slopes, in at most $\rk A$
steps. For uniqueness, in any filtration with semistable decreasing
grades, intersect a submodule with each step; its quotient maps
contractingly into the corresponding grade. Summing degrees bounds its
slope by the first grade's slope, with equality only if all its nonzero
images occur in grades of that slope. Strict decrease therefore places
every maximal-slope submodule in the first term. The first term is $M$,
and induction proves uniqueness. Isometries preserve slopes and ranks,
so uniqueness proves the last assertion.
\end{proof}

\begin{lemma}[Irreducible isometric actions imply semistability]
\label{lem:irreducible-semistable}\uses{lem:hn-filtration}
\source{bost-specht-manuscript:page-0003}\notready
If a finite group acts isometrically on $A$ and its action on $A_K$ is
irreducible, then $A$ is semistable. Independent tensor products and duals
of absolutely irreducible isometric finite-group representations satisfy
the same conclusion.
\end{lemma}
\begin{proof}
The unique maximal destabilizing submodule is preserved by every group
element. Its generic fiber is a nonzero invariant subspace; irreducibility
makes it the full generic fiber. Saturation then makes it all of $A$.

Over an algebraically closed characteristic-zero field, representations
of a finite group are completely reducible: average any projection to an
invariant subspace over the group to obtain an invariant projection.
For an irreducible representation $V$, every equivariant endomorphism
is scalar, since an eigenvalue exists and its eigenspace is invariant.
Consequently, for irreducible $V,W$ of groups $G,H$,
\[
 \operatorname{End}_{G\times H}(V\otimes W)
 =\operatorname{End}_G(V)\otimes\operatorname{End}_H(W)
 \simeq k.
\]
The equality follows by expanding an endomorphism in a basis of one
endomorphism factor and imposing the two independent commuting conditions.
A nontrivial direct-sum decomposition would give a nonscalar projection,
so $V\otimes W$ is irreducible. The annihilator of a proper invariant
subspace of $V^\vee$ would give a proper invariant subspace of $V$;
hence duals are irreducible too. Descent and the first paragraph complete
the proof.
\end{proof}
